\section{Additional Results}
\subsection{Synthetic}
\begin{figure}[ht]
    \centering
     \begin{subfigure}[b]{.32\textwidth}
        \includegraphics[width = \textwidth]{figures/synthetic/statistical/synth_apo_ci_unconfounded.png}
        \caption{APO Functions}
        \label{fig:statistical-uncertainty-apo}   
    \end{subfigure}
    \begin{subfigure}[b]{0.54\textwidth} 
        \includegraphics[width=\textwidth]{figures/synthetic/statistical/synth_capo_ci_unconfounded.png}
        \caption{CAPO Functions}
        \label{fig:statistical-uncertainty-capo}  
    \end{subfigure}
    \caption{Investigating statistical uncertainty using unconfounded synthetic data.}
    \label{fig:statistical-uncertainty}   
\end{figure}

\subsection{Aerosol-Cloud-Climate Effects}

In \Cref{fig:rho} we show how $\Lambda$ can be interpreted as the proportion, $\rho$, of the unexplained range of $\Yt$ attributed to unobserved confounding variables. 
In the left figure, we plot the corresponding bounds for increasing values of $\Lambda$ of the predicted AOD-$\tau$ dose-response curves. 
In the right figure we plot the $\rho$ value for each $\Lambda$ at each value of $\tf$. For the curves reported in \Cref{fig:CODresults_apo}: we find that $\Lambda = 1.1$ leads to $\rho \approx 0.04$, $\Lambda = 1.2$ leads to $\rho \approx 0.07$, and $\Lambda = 1.6$ leads to $\rho \approx 0.15$. 
This shows that when we let just a small amount of the unexplained range of $\Yt$ be attributed to unobserved confounding, the  range  of the predicted APO curves become quite wide. 
If we were to completely relax the no-hidden-confounding assumption, the entire range seen in \Cref{fig:rho} 
Left would be plausible for the APO function. 
This range dwarfs the predicted APO curve. 
These results highlight the importance of reporting such sensitivity analyses.
\begin{figure}
    \centering
    \includegraphics[width=\textwidth]{figures/COD/rho.png}
    \caption{Interpreting $\Lambda$ as a proportion ($\rho$) of the unexplained range of $\Yt$ attributed to unobserved confounding variables.}
    \label{fig:rho}
\end{figure}

In \Cref{fig:cloud-property-dr-curves} we show additional dose response curves for cloud optical thickness ($\tau$), water droplet effective radius ($r_e$), and liquid water path (LWP). 
\begin{figure}[ht]
    \centering
     \begin{subfigure}[b]{.32\textwidth} 
        \includegraphics[width = \textwidth]{figures/COD/dr_tau_attn.png}
        \caption{}
    \end{subfigure}
    \begin{subfigure}[b]{0.32\textwidth}
        \includegraphics[width = \textwidth]{figures/COD/dr_re_attn.png}
        \caption{}
    \end{subfigure}
    \begin{subfigure}[b]{0.32\textwidth}
        \includegraphics[width = \textwidth]{figures/COD/dr_lwp_attn.png}
        \caption{}
    \end{subfigure}
    \caption{Average dose-response curves for other cloud properties. a) Cloud optical depth. b) Water droplet effective radius. c) Liquid water path.}
    \label{fig:cloud-property-dr-curves}
\end{figure}


In \Cref{fig:cloud-property-scatterplots} we show additional scatter plots comparing the neural network and transformer models for cloud optical thickness ($\tau$), water droplet effective radius ($r_e$), and liquid water path (LWP). 
\begin{figure}[ht]
    \centering
     \begin{subfigure}[b]{.32\textwidth} 
        \includegraphics[width = \textwidth]{figures/COD/scatter_tau.png}
        \caption{}
    \end{subfigure}
    \begin{subfigure}[b]{0.32\textwidth}
        \includegraphics[width = \textwidth]{figures/COD/scatter_re.png}
        \caption{}
    \end{subfigure}
    \begin{subfigure}[b]{0.32\textwidth}
        \includegraphics[width = \textwidth]{figures/COD/scatter_lwp.png}
        \caption{}
    \end{subfigure}
    \caption{Comparing transformer to feed-forward feature extractor at predicting cloud properties given covariates and AOD. a) Cloud optical depth. b) Water droplet effective radius. c) Liquid water path. We see a significant improvement in pearson correlation ($R^2$) in each case.}
    \label{fig:cloud-property-scatterplots}
\end{figure}

\subsection{$\omega_{500}$ experiment}

The Overcast models make use of expert knowledge about ACI to select the covariates. 
Ideally, they would include pressure profiles, temperature profiles and supersaturation since these are directly involved in cloud processes and impact the quality of AOD measurements as a proxy for aerosol concentration. 
Unfortunately, they are impossible to retrieve from satellite data, so we rely on meteorological proxies like relative humidity, sea surface temperature, inversion strengths, and vertical motion. 
Relying on these proxies however results in ignorability violations, which must be accounted for in the parameter $\Lambda$ in order to derive appropriate plausible ranges of outcomes.

In the experiment that follows, we are removing a confounding variable from the model, therefore inducing hidden confounding.
The covariate we remove is vertical motion at 500 mb, denoted by $\omega500$. 
This experiment helps us gain some intuition about the influence of the parameter $\Lambda$ and how it relates to the inclusion of confounding variables in the model. 

In \Cref{fig:apo_tr_lrp-vs-lrpw500} we compare the same region with different covariates to identify an appropriate $\Lambda$. 
We fit one model on data from the Pacific (blue) and one model from the Pacific omitting $\omega500$ from the covariates (orange). 
The shaded bounds in blue are the ignorance region for $\Lambda \to 1$ for the Pacific. 
We then find the $\Lambda$ that results in an ignorance interval around the Pacific omitting $\omega500$ that covers the Pacific model prediction. 
From this, we can infer how the parameter $\Lambda$ relates to the inclusion of covariates in the model. We show that we need to set $\Lambda=1.01$ to account for the fact that we are omitting $\omega500$ from our list of covariates.
We also note that the slopes of the dose-response curves are slightly different, with worse predictions when omitting $\omega500$ from the covariates, as expected.

\begin{figure}[htbp]
    \centering
    \begin{subfigure}[b]{0.46\textwidth} 
        \includegraphics[width = \textwidth]{figures/apo_tr_lrp-vs-lrpw500_re.pdf}
        \label{subfig:apo_tr_lrp-vs-lrpw500}
        \caption{Unscaled}
    \end{subfigure}
    \begin{subfigure}[b]{0.46\textwidth} 
        \includegraphics[width = \textwidth]{figures/apo_tr_lrp-vs-lrpw500_lrpw500_re.pdf}
        \label{subfig:apo-scaled_tr_lrp-vs-lrpw500}
        \caption{Scaled by min-max with appropriate $\Lambda$}
    \end{subfigure}
    \caption{Dose-response curves with or without vertical motion at 500 mb ($\omega500$) as a covariate.}
    \label{fig:apo_tr_lrp-vs-lrpw500}
\end{figure}

This work attempts to set a new methodology for setting $\Lambda$ which can be summarised as followed.
Working with two datasets, which vary in only aspect, we train two different models: (i), the control model, and (ii), the experimental model. 
After training both models, we plot the dose-response curves for (i) and (ii) on the same plot. 
We can compare the shape and slope of these curves as well as their uncertainty bounds under the unconfoundedness assumption by plotting the ignorance region for $\Lambda \to 1$ for both models. 
Then, we are interested in setting $\Lambda$ for model (ii) such that the uncertainty bounds cover the entire ignorance region of model (i) under the unconfoundedness assumption. 
For this, we are interested in comparing the slopes and thus min-max scale both curves. 